% day2_seg3_kalstats/meas_noise.tex
\begin{frame}
  \onslide<1->{
	\frametitle{Measurement Noise} 
	\framesubtitle{$v$ and its Covariance Matrix $R$}
}
\begin{itemize}
  \onslide<2->{
\item  $\measnoise{k}$ is the measurement noise sequence;  unlike $w(k)$ it is not simply hypothesized; 
it is physically real and  can be determined empirically
}
  \onslide<3->{
 \item The measurement noise covariance matrix $\measnoisematdef{k}$ indicates the amount of spread or dispersion and internal cross covariance within the measurement itself at any time $k$.
}
  \onslide<4->{
 \item The diagonal elements of $\measnoise{k}$ give the expected variance in each of the individual measurements.
}
\end{itemize}
\end{frame}

\begin{frame}
  \onslide<1->{
	\frametitle{Measurement Noise} 
	\framesubtitle{Characteristics of $R$}
}
\begin{itemize}
  \onslide<2->{
 \item The off-diagonal elements indicate degree of cross covariance between measurement noise random variables 
 associated with each individual measurement;  often 0  
}
  \onslide<3->{
 \item If the noise random variables of individual measurements are uncorrelated, 
 then $R(k)$ will be a diagonal matrix
}
  \onslide<4->{
\item General time series and spectral analysis techniques can be applied to a sequence measurements to determine appropriate values
}
  \onslide<5->{
\item The elements of $R$ are key Kalman tuning parameters
}
\end{itemize}
\end{frame}
